\section{Running-Speed Transfer, Forecast Accuracy, and Closed-Loop Control}
The follow-up uses a fixed public timetable, a 12-vehicle pool and a
0/5/10-s holding menu. It is separate from the budget and supervisor studies.
V412 pairs the capped model-based controller with a frozen within-day
randomized holding rule on the same eight demand days at normal speed and
with the road-speed input reduced by 0.5 m/s. Mean holding and bin shares
pass the registered matching checks in both settings. All four outcome
criteria pass at normal speed. At the slower setting, journey and arrival
CV$^2$ pass while readiness and public-delay nonharm fail
(Table~\ref{tab:running-transfer}).

V414 attempts to forecast actual departure times five minutes before
public release. It replaces the pooled full-link running mean
with the mean at the simulated entry hour when at least three visits at
that hour are strictly completed and available, falling back to the pooled
mean otherwise. Partial links and committed actions do not change.
The two forecasts observe the same capped path. On eight new slow days,
the primary scenario scores \ForecastSlotPairs\ paired supported slots.
The equal-day absolute error falls from \ForecastPooledMAE\ to
\ForecastHourlyMAE\ s, improving on all eight days. Late departures missed
fall from \ForecastPooledMisses\ to \ForecastHourlyMisses, and false alarms
from \ForecastPooledAlarms\ to \ForecastHourlyAlarms. This passes the
forecast screen, not a control test.

V415 executes the hourly and pooled model-based controllers on a different
set of eight slow days. All four registered control tests fail. On one
consumed day examined to explain that failure, two successive launches
advanced by different amounts, widening the interval from 184 to 1,088 s.
Partitioning each actual downstream gap into the difference in dispatch
clocks, accumulated holding and other elapsed time assigns
\FocalLaunchComponent\ of the increase in arrival CV$^2$ to the dispatch
component. Waiting increases by \FocalWaitingLoss\ s/person on that day.
This is covariance accounting on two realized paths, not an estimate of
the causal effect of an individual holding command. All eight days remain
in the closed-loop comparison.
\begin{table}[htbp]
\centering\small
\caption{Running-speed transfer and the separate hourly-predictor deployment test. Entries are paired day means and archived 95\% intervals. Journey is in s/person, readiness in s/trip and public delay in s/day; CV$^2$ is unitless.}
\label{tab:running-transfer}
\setlength{\tabcolsep}{3pt}
\begin{tabular}{@{}lrrrr@{}}\toprule
Comparison & Journey & Readiness & Public delay & Arrival CV$^2$ \\ \midrule
V412 normal & \shortstack[r]{$-3.17$\\{\scriptsize $[-4.64, -1.79]$}} & \shortstack[r]{$-1.84$\\{\scriptsize $[-3.18, -0.35]$}} & \shortstack[r]{$-483.0$\\{\scriptsize $[-832.2, -90.5]$}} & \shortstack[r]{$-0.00623$\\{\scriptsize $[-0.00980, -0.00216]$}} \\[3pt]
V412 slower & \shortstack[r]{$-7.61$\\{\scriptsize $[-13.92, -1.19]$}} & \shortstack[r]{$-4.01$\\{\scriptsize $[-11.05, +3.55]$}} & \shortstack[r]{$-1050.5$\\{\scriptsize $[-2894.3, +930.1]$}} & \shortstack[r]{$-0.04132$\\{\scriptsize $[-0.05238, -0.03061]$}} \\[3pt]
V415 hourly minus pooled & \shortstack[r]{$+1.27$\\{\scriptsize $[-1.25, +5.03]$}} & \shortstack[r]{$-2.19$\\{\scriptsize $[-7.31, +1.46]$}} & \shortstack[r]{$-573.5$\\{\scriptsize $[-1915.4, +381.2]$}} & \shortstack[r]{$+0.00661$\\{\scriptsize $[-0.00376, +0.02403]$}} \\[3pt]
\bottomrule\end{tabular}
\end{table}


V416 also rejects an otherwise admissible holding action if it enlarges a
predicted long departure gap within or at the boundary of a queue of at
least two due trips in one direction. For each modeled pair, the protected
upper bound is the greater of the public headway and the same-state
reference gap, in each existing demand scenario. No release clock or
legal action is added. The comparator is the unprotected hourly controller.
The rule is a local model constraint, not a guarantee on realized gaps.

V417 adds the pooled-plus-guard arm to the same four V416 development days,
reusing the other three arms. Table~\ref{tab:forecast-spacing} reports the
paired differences; the interaction is the hourly guard effect minus the
pooled guard effect. The guard reduces journey on all four days under
hourly prediction, but on only one day under pooled prediction. Both settings
show lower mean actual arrival CV$^2$ with the guard. With both settings
guarded, hourly readiness is lower on all four days, but arrival CV$^2$ is
higher on three and the maximum actual departure gap is longer on all
four. Under this fixed timetable, readiness and public delay summarize the
same lateness at different scales. The arrival-regularity criterion
against unprotected pooled V406 remains the failed V416 development check.
The four-arm comparison is not independent confirmation and does not
justify a policy replacement.
\begin{table}[htbp]
\centering\small
\caption{Forecast-by-spacing development comparison. All contrasts use the same four days. Units match Table~\ref{tab:running-transfer}; Dep. is departure. Intervals are descriptive.}
\label{tab:forecast-spacing}
\setlength{\tabcolsep}{3pt}
\begin{tabular}{@{}lrrrrr@{}}\toprule
Comparison & Journey & Readiness & Public delay & Arrival CV$^2$ & Dep. CV$^2$ \\ \midrule
Hourly, no guard & \shortstack[r]{$+1.20$\\{\scriptsize $[-0.92, +3.40]$}} & \shortstack[r]{$-4.09$\\{\scriptsize $[-6.22, -1.95]$}} & \shortstack[r]{$-1071.0$\\{\scriptsize $[-1630.5, -511.5]$}} & \shortstack[r]{$+0.00527$\\{\scriptsize $[+0.00148, +0.00948]$}} & \shortstack[r]{$+0.00553$\\{\scriptsize $[+0.00173, +0.00989]$}} \\[3pt]
Guard with pooled & \shortstack[r]{$+0.38$\\{\scriptsize $[-0.39, +1.15]$}} & \shortstack[r]{$-0.42$\\{\scriptsize $[-1.31, +0.39]$}} & \shortstack[r]{$-109.0$\\{\scriptsize $[-343.8, +101.5]$}} & \shortstack[r]{$-0.00386$\\{\scriptsize $[-0.00822, +0.00050]$}} & \shortstack[r]{$-0.00426$\\{\scriptsize $[-0.00985, +0.00132]$}} \\[3pt]
Guard with hourly & \shortstack[r]{$-5.15$\\{\scriptsize $[-11.36, -1.98]$}} & \shortstack[r]{$-4.89$\\{\scriptsize $[-12.78, +1.17]$}} & \shortstack[r]{$-1280.2$\\{\scriptsize $[-3348.0, +306.8]$}} & \shortstack[r]{$-0.00417$\\{\scriptsize $[-0.00695, -0.00140]$}} & \shortstack[r]{$-0.00405$\\{\scriptsize $[-0.00669, -0.00043]$}} \\[3pt]
Hourly, both guarded & \shortstack[r]{$-4.34$\\{\scriptsize $[-11.16, +1.01]$}} & \shortstack[r]{$-8.56$\\{\scriptsize $[-14.99, -4.16]$}} & \shortstack[r]{$-2242.2$\\{\scriptsize $[-3927.0, -1089.2]$}} & \shortstack[r]{$+0.00496$\\{\scriptsize $[-0.00338, +0.01330]$}} & \shortstack[r]{$+0.00574$\\{\scriptsize $[-0.00188, +0.01336]$}} \\[3pt]
Interaction & \shortstack[r]{$-5.53$\\{\scriptsize $[-11.63, -1.84]$}} & \shortstack[r]{$-4.47$\\{\scriptsize $[-12.39, +0.81]$}} & \shortstack[r]{$-1171.2$\\{\scriptsize $[-3246.8, +213.0]$}} & \shortstack[r]{$-0.00032$\\{\scriptsize $[-0.00573, +0.00580]$}} & \shortstack[r]{$+0.00021$\\{\scriptsize $[-0.00394, +0.00403]$}} \\[3pt]
V416 vs. pooled & \shortstack[r]{$-3.96$\\{\scriptsize $[-11.34, +1.23]$}} & \shortstack[r]{$-8.97$\\{\scriptsize $[-14.78, -4.71]$}} & \shortstack[r]{$-2351.2$\\{\scriptsize $[-3872.2, -1233.0]$}} & \shortstack[r]{$+0.00110$\\{\scriptsize $[-0.00440, +0.00508]$}} & \shortstack[r]{$+0.00148$\\{\scriptsize $[-0.00288, +0.00449]$}} \\[3pt]
\bottomrule\end{tabular}
\end{table}

